% Standalone manuscript sections; insert after the eventual axis theorem.
% Dependencies: definitions R_N, W_N, F_N, J_N and the proven eventual axis reduction.
% Labels use the requested eul: namespace.  All constants are reintroduced.

\section{A second asymptotic scale from the gamma transition region}
\label{eul:sec:second-scale}
The elementary Taylor envelope gives a power remainder and proves the
matching leading constant.  Its positive remainder has a more precise
asymptotic.  The relevant gamma mass is centered at the transition
$t=\tfrac12\log N$, rather than at either of the two moment saddles
that produce $2^{-b}$ and $N3^{-b}$.

Retain
\[
 d=\log(3/2),\quad \lambda=d^{-1},\quad p=\lambda\log2,
 \quad q=\frac{\log3}{\log2},\quad
 B_0=\lambda\log q,\quad
 \widehat b_N=\lambda\log N+B_0.
\]
Since $q=1+1/p$, the leading amplitude has the equivalent form
\begin{equation}\label{eul:eq:c-simplified}
 c=(1-q^{-1})q^{-p}=\frac{p^p}{(p+1)^{p+1}}.
\end{equation}
Introduce the additional constants
\begin{align}
 \alpha&=\lambda-\tfrac12,\qquad \rho=2\lambda-1=2\alpha,
 &\mu&=\tfrac12-\lambda+\lambda\log(2\lambda),\label{eul:eq:mu}\\
 D_0&=\frac{\Gamma(-\alpha)}2
       \sqrt{\frac{2\lambda}{\pi}}(2\lambda)^{-B_0},
 &K&=(\log2)dq^{-p},\label{eul:eq:second-amplitudes}\\
 D_1&=\frac{\log(2\lambda)D_0}{K}.
 \nonumber
\end{align}
Here $1<\alpha<2$, so $\Gamma(-\alpha)>0$.
The exponents obey
\begin{equation}\label{eul:eq:exponent-order}
 p<\mu<2p-1.
\end{equation}
For orientation, numerical evaluations are
\[
 \mu=1.9695904144700607703\ldots,\quad
 \mu-p=0.2600791231186059934\ldots,
\]
\[
 D_0=1.5677702007103631657\ldots,\qquad
 D_1=19.562981709827216349\ldots.
\]
The exact formulas, rather than these decimals, enter the proofs.

\begin{theorem}[The next term and the direction of the maximizing shift]
\label{eul:thm:second-term}
Let $b_N$ be the unique axis maximizer, which also gives the global
constant $C_N$ for every sufficiently large $N$.  Then
\begin{align}
 C_N&=1+cN^{-p}
       +D_0\frac{N^{-\mu}}{\sqrt{\log N}}
       +o\!\left(\frac{N^{-\mu}}{\sqrt{\log N}}\right),
       \label{eul:eq:CN-second}\\
 b_N&=\widehat b_N
       -D_1\frac{N^{-(\mu-p)}}{\sqrt{\log N}}
       +o\!\left(\frac{N^{-(\mu-p)}}{\sqrt{\log N}}\right).
       \label{eul:eq:bN-second}
\end{align}
In particular the true maximizing inner order is eventually smaller
than $\widehat b_N$, and the leading approximation
$1+cN^{-p}$ eventually underestimates the sharp constant.
\end{theorem}

\subsection{The exact remainder and its transition limit}
Put
\[
 G_N(z)=\frac{(1-z)^{N-1}}{1+z},\qquad
 r_N(z)=G_N(z)-1+Nz.
\]
The earlier exact axis expansion reads
\begin{equation}\label{eul:eq:axis-second-start}
 R_N(0,b)=1+2^{-b}-N3^{-b}-N4^{-b}+\varepsilon_N(b),
\end{equation}
where
\begin{equation}\label{eul:eq:epsilon-integral}
 \varepsilon_N(b)=\frac1{\Gamma(b)}\int_0^\infty
 t^{b-1}e^{-t}(1+e^{-t})r_N(e^{-2t})\,dt.
\end{equation}
The nonnegative Taylor remainder satisfies both bounds
\begin{equation}\label{eul:eq:two-remainder-bounds}
 0\leq r_N(z)\leq\min\{Nz,A_Nz^2\},\qquad
 A_N=(N^2-N+2)/2.
\end{equation}
The quadratic bound is the previous Taylor theorem.  The linear bound
uses $0\leq G_N(z)\leq1$.

\begin{lemma}[Uniform transition asymptotics with two derivatives]
\label{eul:lem:transition-C2}
Let $\ell=\log N$, let $B$ range over any fixed compact real interval,
and set $b=\lambda\ell+B$.  Uniformly in $B$, for $j=0,1,2$,
\begin{equation}\label{eul:eq:epsilon-C2}
 \frac{d^j}{db^j}\varepsilon_N(b)
 =\frac{\Gamma(-\alpha)}2\sqrt{\frac{2\lambda}{\pi}}
  (2\lambda)^{-B}\{-\log(2\lambda)\}^{j}
  \frac{N^{-\mu}}{\sqrt\ell}
  +o\!\left(\frac{N^{-\mu}}{\sqrt\ell}\right).
\end{equation}
\end{lemma}
\begin{proof}
Set $L=\ell/2$ and
\[
 P_N(B)=\frac{e^{-L}L^{b-1}}{\Gamma(b)}.
\]
With $t=L+y$, equation \eqref{eul:eq:epsilon-integral} becomes
\begin{equation}\label{eul:eq:transition-rescaling}
 \frac{\varepsilon_N(b)}{P_N(B)}
 =\int_{-L}^\infty D_{N,B}(y)
       (1+N^{-1/2}e^{-y})r_N(N^{-1}e^{-2y})\,dy,
\end{equation}
where
\[
 D_{N,B}(y)=
  \exp\{(b-1)\log(1+y/L)-y\}.
\]
For each fixed $y$, uniformly for bounded $B$,
\[
 D_{N,B}(y)\longrightarrow e^{\rho y},\qquad
 r_N(N^{-1}e^{-2y})\longrightarrow
              e^{-e^{-2y}}-1+e^{-2y}.
\]
The integrands vanish below $-L$.

Choose a fixed $\epsilon>0$ so small that
$2<\rho-3\epsilon<\rho+3\epsilon<4$.
The inequality $\log(1+u)\leq u$ and boundedness of $B$ give, for
all sufficiently large $N$,
\[
 D_{N,B}(y)\leq
 \begin{cases}
   e^{(\rho+\epsilon)y},&y\geq0,\\
   e^{(\rho-\epsilon)y},&-L<y<0.
 \end{cases}
\]
Also $1+N^{-1/2}e^{-y}\leq2$ on this domain.  Equation
\eqref{eul:eq:two-remainder-bounds} bounds the remaining factor by
$\min(e^{-2y},C e^{-4y})$, with a fixed $C$.
The resulting dominating function is integrable at both infinities:
use the quartic bound for $y\geq0$ and the quadratic-in-$e^{-y}$
bound for $y<0$.  Dominated convergence therefore gives
\begin{equation}\label{eul:eq:Mellin-limit}
 \frac{\varepsilon_N(b)}{P_N(B)}
 \longrightarrow
 I_\rho:=\int_{-\infty}^{\infty}
      e^{\rho y}(e^{-e^{-2y}}-1+e^{-2y})\,dy
 =\frac12\Gamma(-\alpha).
\end{equation}
To verify the last evaluation, set $x=e^{-2y}$:
\[
 I_\rho=\frac12\int_0^\infty
          x^{-\alpha-1}(e^{-x}-1+x)\,dx.
\]
Two integrations by parts have zero boundary terms because
$1<\alpha<2$, and give
$I_\rho=\Gamma(2-\alpha)/(2\alpha(\alpha-1))
       =\Gamma(-\alpha)/2$.

For clarity, the same domination controls parameter derivatives;
this is not an inference from differentiating an asymptotic formula.
The first two derivatives of the gamma density insert the factors
\[
 S_{N,B}(y)=\log(L+y)-\psi(b),\qquad
 S_{N,B}(y)^2-\psi_1(b).
\]
At each fixed $y$ these tend uniformly to $-\log(2\lambda)$ and
$\log^2(2\lambda)$, respectively.
For $y\geq0$, their absolute values are bounded by a fixed multiple
of $1+y^2$, which is integrable against the preceding exponential
majorant after a harmless adjustment of $\epsilon$.
For negative $y$ write $u=1+y/L\in(0,1)$.
For $j=1,2$ and large $L$,
\[
 u^{b-1}|\log u|^j\leq C_j u^{b-1-\epsilon L},
\]
because $\sup_{v\geq0}v^j e^{-\epsilon L v}$ is uniformly bounded.
Replacing $b-1$ by $b-1-\epsilon L$ in the preceding exponential
estimate still leaves an exponent strictly greater than two on the
negative half-line.  The terms $\log L-\psi(b)$ and $\psi_1(b)$
are uniformly bounded.  This supplies a common integrable majorant
for the differentiated integrands as well.

Convergence is uniform for bounded $B$: the pointwise convergences
are uniform on each compact $y$ interval and the same integrable
majorants apply to every $B$ in the prescribed compact interval.
Finally, uniform Stirling asymptotics \cite{DLMFGamma} at $b=\lambda\ell+B$ give
\begin{equation}\label{eul:eq:gamma-window-density}
 P_N(B)=\sqrt{\frac{2\lambda}{\pi}}(2\lambda)^{-B}
           \frac{N^{-\mu}}{\sqrt\ell}\{1+O(\ell^{-1})\}.
\end{equation}
Combining this with the three dominated-convergence limits proves
\eqref{eul:eq:epsilon-C2}.
\end{proof}

\subsection{Optimization at the second scale}
\begin{proof}[Proof of Theorem~\ref{eul:thm:second-term}]
We first record exact inequalities between the exponents, to avoid
using numerical approximations as premises.  The elementary bounds
$2/5<\log(3/2)<5/12$ give $12/5<\lambda<5/2$.
The function
$x\mapsto\tfrac12-x+x\log x$ is increasing for $x>1$, and is
positive at $x=12/5$; for the latter assertion,
$\log(12/5)>2(12/5-1)/(12/5+1)=14/17$ suffices.
It follows that $\mu-p>0$.
Also
\[
 \mu-(2p-1)=\tfrac32-\lambda+\lambda\log(\lambda/2)<0:
\]
the expression is increasing on $2<\lambda<5/2$, and at $5/2$ it
is $-1+(5/2)\log(5/4)<-3/8$.
This proves \eqref{eul:eq:exponent-order}.

Define the two scales
\[
 h_N=\frac{N^{-\mu}}{\sqrt\ell},\qquad
 t_N=N^ph_N=\frac{N^{-(\mu-p)}}{\sqrt\ell}.
\]
For bounded $s$, set
\[
 f(s)=q^{-p}2^{-s}-q^{-(p+1)}3^{-s},\qquad
 g(s)=D_0(2\lambda)^{-s}.
\]
Equation \eqref{eul:eq:axis-second-start}, Lemma
\ref{eul:lem:transition-C2}, and $\mu<2p-1$ give the expansion
\begin{equation}\label{eul:eq:two-scale-C2}
 R_N(0,\widehat b_N+s)
   =1+N^{-p}f(s)+h_Ng(s)+o(h_N),
\end{equation}
uniformly with its first two derivatives in $s$ on every compact
interval.  In fact the omitted term $N4^{-b}$ and both of its
derivatives are $O(N^{-(2p-1)})=o(h_N)$.

The already proved localization gives $s_N=b_N-\widehat b_N\to0$.
Differentiating the integral through the justified
\eqref{eul:eq:two-scale-C2}, the stationary equation is
\[
 0=f'(s_N)+t_Ng'(s_N)+o(t_N).
\]
Since $f'(0)=0$ and $f''(0)=-K<0$, this first implies
$s_N=O(t_N)$, and then
\[
 -Ks_N-t_N\log(2\lambda)D_0+o(t_N)=0.
\]
This is \eqref{eul:eq:bN-second}.
Finally $f(s_N)=c+O(t_N^2)$ and $g(s_N)=D_0+O(t_N)$.
The displacement changes the height by
$O(N^{-p}t_N^2)=O(h_Nt_N)=o(h_N)$.
Substituting into \eqref{eul:eq:two-scale-C2} proves
\eqref{eul:eq:CN-second}.
\end{proof}

\subsection{The first inverse-logarithmic correction}
The convergence of the second-scale amplitude is particularly slow
because $\alpha$ is close to two.  The first correction describes
this effect explicitly.  Let $\psi$ and $\psi_1$ denote the digamma
and trigamma functions, continued to $-\alpha$, which is not a pole.
Put
\begin{equation}\label{eul:eq:T-B}
 T(B)=-(B-1)\psi(-\alpha)
 -\frac\lambda2\{\psi_1(-\alpha)+\psi(-\alpha)^2\}
 -\frac{B^2-B+1/6}{2\lambda}.
\end{equation}

\begin{proposition}[An explicit logarithmic refinement]
\label{eul:prop:logarithmic-refinement}
Uniformly for bounded $B$,
\begin{equation}\label{eul:eq:epsilon-log-refinement}
 \varepsilon_N(\lambda\ell+B)
 =\frac{\Gamma(-\alpha)}2\sqrt{\frac{2\lambda}{\pi}}
   (2\lambda)^{-B}\frac{N^{-\mu}}{\sqrt\ell}
   \left\{1+\frac{T(B)}\ell+O(\ell^{-2})\right\}.
\end{equation}
Consequently the global constant has the refinement
\begin{equation}\label{eul:eq:CN-log-refinement}
 C_N=1+cN^{-p}
  +D_0\frac{N^{-\mu}}{\sqrt\ell}
     \left\{1+\frac{T(B_0)}\ell+O(\ell^{-2})\right\}.
\end{equation}
\end{proposition}
\begin{proof}
All moments in $y$ of the limiting integrand in
\eqref{eul:eq:Mellin-limit} are finite, by the same two exponential
tail bounds.  On a central interval $|y|\leq M\log\ell$, Taylor
expansion gives
\[
 D_{N,B}(y)=e^{\rho y}
 \left\{1+\frac{(B-1)y-\lambda y^2}{L}
     +O\!\left(\frac{1+|y|^6}{L^2}\right)\right\},
\]
uniformly for bounded $B$; the harmless constant in the remainder
can be chosen uniformly for all sufficiently large $\ell$.
On the same interval the finite-$N$ kernel can be replaced by its
limit with an error smaller than every fixed negative power of
$\ell$.  Indeed $e^{-2y}=O(\ell^{2M})$,
$N^{-1}e^{-2y}\to0$ faster than every power of $1/\ell$,
and expanding $\log G_N(N^{-1}e^{-2y})$ around zero gives this
assertion; the factor $1+N^{-1/2}e^{-y}$ has the same property.
For each fixed target power of $1/\ell$, choose $M$ large enough
that the two uniform exponential tail bounds from Lemma
\ref{eul:lem:transition-C2} make the integrals outside this central
interval smaller than that power.  Integrating the displayed Taylor
expansion, and extending its polynomial moments to the entire real
line, therefore gives
\[
 \frac{\varepsilon_N(\lambda\ell+B)}{P_N(B)}
 =I_\rho+\frac{(B-1)I_\rho'-\lambda I_\rho''}{L}
       +O(L^{-2}).
\]
Here primes differentiate the explicit function
$I_\rho=\Gamma(-\rho/2)/2$ in $\rho$; hence
\[
 \frac{I_\rho'}{I_\rho}=-\tfrac12\psi(-\alpha),\qquad
 \frac{I_\rho''}{I_\rho}=
       \tfrac14\{\psi_1(-\alpha)+\psi(-\alpha)^2\}.
\]
The uniform Stirling expansion one order beyond
\eqref{eul:eq:gamma-window-density} is
\[
 P_N(B)=\sqrt{\frac{2\lambda}{\pi}}(2\lambda)^{-B}
 \frac{N^{-\mu}}{\sqrt\ell}
 \left\{1-\frac{B^2-B+1/6}{2\lambda\ell}
              +O(\ell^{-2})\right\}.
\]
Multiplication and $L=\ell/2$ give
\eqref{eul:eq:epsilon-log-refinement}.
For \eqref{eul:eq:CN-log-refinement}, the optimizing shift changes
the height by $O(h_Nt_N)=o(h_N\ell^{-2})$, and
$N4^{-\widehat b_N}=o(h_N\ell^{-2})$ by the strict exponent
inequality.  Thus one can evaluate the refined remainder at $B=B_0$.
\end{proof}

For orientation, $T(B_0)=-2098.4287350903144\ldots$.
This large coefficient explains why numerical indices such as $10^8$
already illustrate the leading law well but do not yet closely
approximate the second-scale amplitude $D_0$.  The theorem is an
asymptotic statement as $N$ tends to infinity; no practical onset
index for the second expansion is asserted.

\section{Strict log-convexity of the axis envelope}
\label{eul:sec:axis-logconvex}
For real $t\geq1$, define the continuous interpolation
\[
 \mathcal A(t)=\sup_{b>0}\frac1{\Gamma(b)}\int_0^\infty
 u^{b-1}e^{-u}
 \frac{(1+e^{-u})(1-e^{-2u})^{t-1}}{1+e^{-2u}}\,du.
\]
At integer indices, $\mathcal A(N)$ is the optimized axis constant.
For every real $t\geq1$, the supremum is a finite maximum greater
than one.  To verify existence, note the endpoint limits in $b$:
the limit at infinity is one, while the limit at zero is zero if
$t>1$ and one if $t=1$.  Also the kernel has expansion
$1+v-tv^2+O_t(v^3)$ at $v=0$.  Its gamma average consequently equals
$1+2^{-b}-t3^{-b}+O_t(4^{-b})$ as $b\to\infty$, and is greater
than one for sufficiently large $b$.
The last big-oh bound follows by extending the bounded Taylor
remainder divided by $v^3$ continuously near zero and bounding it
on the rest of $[0,1]$.

\begin{theorem}[Strict log-convex envelope]
\label{eul:thm:axis-logconvex}
For distinct real $s,t\geq1$ and $0<\theta<1$,
\[
 \mathcal A(\theta s+(1-\theta)t)
     <\mathcal A(s)^\theta\mathcal A(t)^{1-\theta}.
\]
In particular, at every integer $N\geq2$,
\[
 \mathcal A(N)^2<\mathcal A(N-1)\mathcal A(N+1).
\]
The globally optimized Euler constants satisfy
$C_N^2<C_{N-1}C_{N+1}$ for every sufficiently large $N$.
\end{theorem}
\begin{proof}
For fixed $b$, its axis integral is a positive moment of
$(1-e^{-2u})^{t-1}$.  Strict H\"older inequality gives strict
log-convexity in $t$, since $\log(1-e^{-2u})$ is not constant
under its positive gamma-weighted density.
Choose a finite maximizing $b$ for the intermediate argument
$\theta s+(1-\theta)t$.  Apply strict H\"older at that one value
of $b$, and bound its two endpoint integrals by their respective
maxima.  This yields the asserted strict inequality; it avoids
the false general inference that an arbitrary supremum preserves
\emph{strict} convexity without an attained maximizer.
The discrete conclusions follow by $\theta=1/2$ and the eventual
axis reduction.
\end{proof}
